% Generated from tex/chapters/ch04/07_実務上の考慮.tex for the SWEBOK structure tree
\subsubsection{クロスプラットフォーム開発と移行}

モバイルアプリケーションなどは、特定のプラットフォームに強く依存することが多い。たとえば、iOS と Android では、開発環境、API、利用できる機能、利用者体験が異なる。各プラットフォーム向けに別々に開発すると、費用と時間が増え、実装差も生まれやすい。

クロスプラットフォーム開発は、共通の言語やフレームワークで作り、複数のプラットフォームに出力する方法である。方式としては、共通言語から各プラットフォーム向けのネイティブアプリを生成する方法と、Web 技術で作った部分をネイティブのコンテナで包むハイブリッド方式がある。

既存アプリケーションを別のプラットフォームへ移す場合は、移行が必要になる。移行では、プログラミング言語の違い、API の違い、データ形式、テスト環境、利用者体験の差を考慮する必要がある。
